\documentclass{subfiles}
\begin{document}
    Let us consider an instance of the TSP in which the cost function is such that the 
        triangular inequality holds; that is, given the edges \((u, w), (w, v)\) and \(u, v)\)
        it holds 
        \[
            c(u, w) + c(w, v) \le c(u, v).
        \]
        We show that the following algorithm solves TSP and 
        that is a 2-approximation algorithm.

    The idea is the following: let \(G = (V, E)\) be a complete graph,
        and let \(T\) be the spanning tree with root \(r \in V\). 
        Then, if \(H\) is the list of vertexes encountered in the pre-order view of \(T\),
        we compute a tour by simply considering the edges between any two consectutive 
        nodes in \(H\). We give the pseudocode in \autoref{proc:approx-tsp}.

        \subfile{../TikZ/Procedure - Approximate TSP}

    \begin{theorem*}
        \autoref{proc:approx-tsp} is a polynomial 2-approximate algorithm for the 
            traveling-salesman problem.
    \end{theorem*}

    \begin{proof*}
        Let \(H^{*}\) be an optimal tour for the given graph.
            From \(H^{*}\), we obtain a spanning tree \(T\) by deliting edges from the tour.
            Hance the algorithm provides a lower bound for an optimal tour:
            \[
                c(T) \le c(H^{*}).
            \]
            If we denote by \(W\) a full walk on \(T\),
                and note that such walk traverses each node exactly twice,
                one has that 
            \[
                c(W) = 2c(T).
            \]
            Combining the two equations, it follows:
            \[
                c(W) \le 2c(H^{*}).
            \]

        Now, \(W\) may not be a tour. Thankfully we can employ the triangular inequality to 
            to delete a visit to a node in \(W\). Repeating the procedure,
            we can remove all but the first visit to each vertex.
            This ordering is the same as the pre-order visit of \(T\).
            Let \(H\) be such a ordering. Since \(H\) is obtained from \(W\),
            it follows
            \[
                c(H) \le c(W),
            \]
            hence, we get 
            \[
                c(H) \le 2c(H^{*})
            \]
    \end{proof*}
\end{document}
